\documentclass[10pt,a4paper]{amsart}
\usepackage{amsmath,amssymb,amsthm,graphicx}
% Retained preamble for the Stacks Project chapter snapshot in this directory.
% The chapter's own preamble.tex is a copy of the shared Stacks Project preamble; the block
% below is sliced MECHANICALLY, line for line, from it. Nothing was added, removed or
% redefined; the wrapper class, the xr-hyper/hyperref wiring and the externaldocument list
% are omitted because the assembler supplies the document class.
%
% lines taken from preamble.tex: 3, 5, 9, 11, 12, 156, 157, 158, ... (51 lines total)
%
% package coverage restored from the chapter's own preamble (2026-10-13): xy, xy-2cell

\usepackage[all]{xy}
\xyoption{2cell}
\UseAllTwocells

\newenvironment{reference}{\comment}{\endcomment}
\newenvironment{slogan}{\comment}{\endcomment}
\newenvironment{history}{\comment}{\endcomment}
\theoremstyle{plain}
\newtheorem{theorem}[subsection]{Theorem}
\newtheorem{proposition}[subsection]{Proposition}
\newtheorem{lemma}[subsection]{Lemma}
\theoremstyle{definition}
\newtheorem{definition}[subsection]{Definition}
\newtheorem{example}[subsection]{Example}
\newtheorem{exercise}[subsection]{Exercise}
\newtheorem{situation}[subsection]{Situation}
\theoremstyle{remark}
\newtheorem{remark}[subsection]{Remark}
\newtheorem{remarks}[subsection]{Remarks}
\numberwithin{equation}{subsection}
\def\lim{\mathop{\mathrm{lim}}\nolimits}
\def\colim{\mathop{\mathrm{colim}}\nolimits}
\def\Spec{\mathop{\mathrm{Spec}}}
\def\Hom{\mathop{\mathrm{Hom}}\nolimits}
\def\Ext{\mathop{\mathrm{Ext}}\nolimits}
\def\SheafHom{\mathop{\mathcal{H}\!\mathit{om}}\nolimits}
\def\SheafExt{\mathop{\mathcal{E}\!\mathit{xt}}\nolimits}
\def\Sch{\mathit{Sch}}
\def\Mor{\mathop{\mathrm{Mor}}\nolimits}
\def\Ob{\mathop{\mathrm{Ob}}\nolimits}
\def\Sh{\mathop{\mathit{Sh}}\nolimits}
\def\NL{\mathop{N\!L}\nolimits}
\def\CH{\mathop{\mathrm{CH}}\nolimits}
\def\proetale{{pro\text{-}\acute{e}tale}}
\def\etale{{\acute{e}tale}}
\def\QCoh{\mathit{QCoh}}
\def\Ker{\mathop{\mathrm{Ker}}}
\def\Im{\mathop{\mathrm{Im}}}
\def\Coker{\mathop{\mathrm{Coker}}}
\def\Coim{\mathop{\mathrm{Coim}}}
\DeclareMathSymbol{\boxtimes}{\mathbin}{AMSa}{"02}
\def\QCohstack{\mathcal{QC}\!\mathit{oh}}
\def\Cohstack{\mathcal{C}\!\mathit{oh}}
\def\Spacesstack{\mathcal{S}\!\mathit{paces}}
\def\Quotfunctor{\mathrm{Quot}}
\def\Hilbfunctor{\mathrm{Hilb}}
\def\Curvesstack{\mathcal{C}\!\mathit{urves}}
\def\Polarizedstack{\mathcal{P}\!\mathit{olarized}}
\def\Complexesstack{\mathcal{C}\!\mathit{omplexes}}
\def\Pic{\mathop{\mathrm{Pic}}\nolimits}
\def\Picardstack{\mathcal{P}\!\mathit{ic}}
\def\Picardfunctor{\mathrm{Pic}}
\def\Deformationcategory{\mathcal{D}\!\mathit{ef}}
\usepackage{cleveref}
\crefname{originallemma}{original lemma}{original lemmas}
\Crefname{originallemma}{Original lemma}{Original lemmas}
\makeatletter
\newlabel{lemma-rings-TI}{{0DY3}{source}{}{sourcepointer-1}{}}
\newlabel{lemma-rings-TI@cref}{{[originallemma][0][]0DY3}{[1][0][]source}{}{}{}}
\makeatother
\title{Stacks, Examples of Deformations: lemma-power-series-rings-TI}
\author{The Stacks Project Authors}
\date{stacks-project:examples-defos.tex:local-snapshot-20261012; GNU FDL 1.2 or later}
\begin{document}
\maketitle
\noindent\textit{Editorial note.} Statement span is lines 1341-1354 of sources/stacks-examples-defos/pinned-original-source.tex; proof span is lines 1356-1411, the complete author proof (54 body lines).\par\medskip
\noindent\textit{Editorial note.} Navigation stubs: 1 cited-but-not-included labels are supplied as non-load-bearing pointers with their pinned tags where the snapshot tag index has them.\par\medskip
\noindent\textit{Editorial note.} extra\_packages is empty on purpose: the retained preamble already loads xy with options and a second declaration clashes in the assembler wrapper.\par\medskip
\noindent\textit{Editorial note.} Text is retained verbatim, including the author's own typos and placeholders.\par\medskip
\noindent\textit{Editorial note.} Queue index 8 of queue\_examples-defos.json (pb=2052, stmt 1341-1354, proof 1356-1411).\par\medskip
\noindent\textit{Allowed standard prerequisites (editorial).}\begin{itemize}
\item Explicit prerequisite boundary. The cited inputs are prerequisites; the argument in the retained span is not replaced by any cited result.
\end{itemize}\medskip
\ifdefined\hypertarget\hypertarget{sourcepointer-1}{}\fi
\noindent\textit{Source reference.} Original lemma 0DY3: sources/stacks-examples-defos/pinned-original-source.tex, source lines 1072--1079. Cited by the retained statement or proof but not part of the collected spans; supplied only as a navigation pointer. The argument does not depend on its content.\par\medskip
\begin{lemma}
\label{lemma-power-series-rings-TI}
In Lemma \ref{lemma-rings-TI} if $P = k[[x_1, \ldots, x_n]]/(f)$
for some nonzero $f \in (x_1, \ldots, x_n)^2$, then
\begin{enumerate}
\item $\text{Inf}(\Deformationcategory_P)$ is finite dimensional
if and only if $n = 1$, and
\item $T\Deformationcategory_P$ is finite dimensional if
$$
\sqrt{(f, \partial f/\partial x_1, \ldots,  \partial f/\partial x_n)} =
(x_1, \ldots, x_n)
$$
\end{enumerate}
\end{lemma}

\begin{proof}
Proof of (1). Consider the derivations $\partial/\partial x_i$ of
$k[[x_1, \ldots, x_n]]$ over $k$. Write $f_i = \partial f/\partial x_i$.
The derivation
$$
\theta = \sum h_i \partial/\partial x_i
$$
of $k[[x_1, \ldots, x_n]]$
induces a derivation of $P = k[[x_1, \ldots, x_n]]/(f)$
if and only if
$\sum h_i f_i \in (f)$. Moreover, the induced derivation of $P$
is zero if and only if $h_i \in (f)$ for $i = 1, \ldots, n$.
Thus we find
$$
\Ker((f_1, \ldots, f_n) : P^{\oplus n} \longrightarrow P) \subset
\text{Der}_k(P, P)
$$
The left hand side is a finite dimensional $k$-vector space only if
$n = 1$; we omit the proof. We also leave it to the reader to see
that the right hand side has finite dimension if $n = 1$.
This proves (1).

\medskip\noindent
Proof of (2). Let $Q$ be a flat deformation of $P$ over $k[\epsilon]$
as in the proof of Lemma \ref{lemma-rings-TI}. Choose lifts $q_i \in Q$
of the image of $x_i$ in $P$. Then $Q$ is a complete local ring
with maximal ideal generated by $q_1, \ldots, q_n$ and $\epsilon$
(small argument omitted). Thus we get a surjection
$$
k[\epsilon][[x_1, \ldots, x_n]] \longrightarrow Q,\quad
x_i \longmapsto q_i
$$
Choose an element of the form
$f + \epsilon g \in k[\epsilon][[x_1, \ldots, x_n]]$
mapping to zero in $Q$. Observe that $g$ is well defined modulo $(f)$.
Since $Q$ is flat over $k[\epsilon]$ we get
$$
Q = k[\epsilon][[x_1, \ldots, x_n]]/(f + \epsilon g)
$$
Finally, if we changing the choice of $q_i$ amounts to
changing the coordinates $x_i$ into $x_i + \epsilon h_i$
for some $h_i \in k[[x_1, \ldots, x_n]]$. Then
$f + \epsilon g$ changes into $f + \epsilon (g + \sum h_i f_i)$
where $f_i = \partial f/\partial x_i$. Thus we see that the
isomorphism class of the deformation $Q$ is determined
by an element of
$$
k[[x_1, \ldots, x_n]]/
(f, \partial f/\partial x_1, \ldots,  \partial f/\partial x_n)
$$
This has finite dimension over $k$ if and only if
its support is the closed point of $k[[x_1, \ldots, x_n]]$
if and only if
$\sqrt{(f, \partial f/\partial x_1, \ldots,  \partial f/\partial x_n)} =
(x_1, \ldots, x_n)$.
\end{proof}

\end{document}
